\subsection*{Problema 08}

\textbf{Enunciado:}
Calcular la integral indefinida:
\begin{align*}
\int \dfrac{dx}{\sqrt{x^2 + x + 1}}
\end{align*}

\textbf{Solución:}
Para resolver la integral, primero completamos el cuadrado en el trinomio del denominador:
\begin{align*}
x^2 + x + 1 &= \left(x + \dfrac{1}{2}\right)^2 \\
&\quad + 1 - \left(\dfrac{1}{2}\right)^2
\end{align*}

Simplificando los términos numéricos, obtenemos la expresión equivalente:
\begin{align*}
x^2 + x + 1 = \left(x + \dfrac{1}{2}\right)^2 + \dfrac{3}{4}
\end{align*}

Sustituimos esto en la integral original:
\begin{align*}
\int \dfrac{dx}{\sqrt{\left(x + \dfrac{1}{2}\right)^2 + \dfrac{3}{4}}}
\end{align*}

Utilizamos la fórmula de integración directa para funciones irracionales:
\begin{align*}
\int \dfrac{du}{\sqrt{u^2 + a^2}} = \ln\left|u + \sqrt{u^2 + a^2}\right| + C
\end{align*}

En nuestro caso, identificamos las variables:
\begin{align*}
u &= x + \dfrac{1}{2}, \quad dx = du \\
a^2 &= \dfrac{3}{4} \implies a = \dfrac{\sqrt{3}}{2}
\end{align*}

Aplicando la fórmula de integración, el resultado preliminar es:
\begin{align*}
\ln\left|x + \dfrac{1}{2} + \sqrt{x^2 + x + 1}\right| + C
\end{align*}

Por lo tanto, la solución general de la integral es:
\begin{align*}
\int \dfrac{dx}{\sqrt{x^2 + x + 1}} &= \ln\left|x + \dfrac{1}{2} \right. \\
&\quad \left. + \sqrt{x^2 + x + 1}\right| + C
\end{align*}

\begin{figure}[H]
	\centering
	\includegraphics[width=\columnwidth]{semana_05/graficas/SEM05_P08_grafica.png}
	\caption{Representación gráfica de $f(x)$ y su antiderivada $F(x)$ evaluada para la constante de integración $C = 0$ (Problema 08).}
	\label{fig:problema08}
\end{figure}
